% =====================================================================
\chapter{Συζήτηση}
\label{ch:discussion}
% =====================================================================

\section{Ικανοποίηση των στόχων}
\label{sec:disc-goals}

Ο Πίνακας~\ref{tab:goals} αντιστοιχίζει τους στόχους του Κεφαλαίου~\ref{ch:intro}
με τους μηχανισμούς και την τεκμηρίωση που τους ικανοποιούν.

\begin{table}[htbp]
\centering
\small
\caption{Αντιστοίχιση στόχων, μηχανισμών και τεκμηρίωσης.}
\label{tab:goals}
\begin{tabular}{@{}p{3.4cm}p{5.6cm}p{4.6cm}@{}}
\toprule
\textbf{Στόχος} & \textbf{Μηχανισμός} & \textbf{Τεκμηρίωση}\\
\midrule
Σ1 Προεπιλεγμένη άρνηση & \code{authorized\_default=0}, κράτηση χωρίς πράκτορα & \code{test\_gate}, \code{test\_daemon}\\
Σ2 Τεκμηριωμένη απόφαση & αναφορά με ευρήματα 3 επιπέδων & §\ref{sec:example-badusb}\\
Σ3 Ανίχνευση αντιφάσεων & 21 κανόνες σταδίου 2 & Πίνακας~\ref{tab:scenarios}\\
Σ4 Παρατήρηση συμπεριφοράς & \code{EVIOCGRAB}, \code{unprompted-typing} & Πίνακας~\ref{tab:behaviour}\\
Σ5 Μνήμη & ιστορικό, \code{descriptor-drift} & §\ref{sec:example-drift}\\
Σ6 Ελάχιστο προνόμιο & \code{--privsep}, επαληθευμένη απόρριψη & \code{test\_privsep}, \code{test\_boundary}, κάλυψη 95\%\\
Σ7 Αποφυγή κλειδώματος & 4 διαδρομές επαναφοράς, \en{watchdog}, αρχείο πανικού & \code{test\_gate}, \code{test\_service}\\
Σ8 Επαληθευσιμότητα & 1413 έλεγχοι χωρίς υλικό, \en{testbed} & Κεφάλαιο~\ref{ch:evaluation}\\
\bottomrule
\end{tabular}
\end{table}

\section{Περιορισμοί}
\label{sec:limitations}

Η εργασία καταγράφει ρητά τις αδυναμίες του συστήματος, ώστε ο χρήστης να μην
το θεωρήσει πανάκεια.

\begin{enumerate}
  \item \textbf{Το κενό έκθεσης του σταδίου 3.} Η συσκευή ενεργοποιείται πριν
        δεσμευτεί η είσοδός της. Πλήκτρα που στέλνονται σε αυτό το διάστημα
        (41–85\,\en{ms} σε πραγματικό υλικό) φτάνουν στη συνεδρία. Η πειραματική
        επιλογή \code{--close-race-window} (καθυστερημένη πρόσδεση οδηγού) δεν
        κλείνει πλήρως το κενό. Για πλήρη απομόνωση, ο χρήστης μπορεί να
        παραλείψει τα στάδια 3–4.
  \item \textbf{Τέλεια αντίγραφα.} Ένα καλώδιο \en{O.MG} που δηλώνει ακριβώς ό,τι
        ένα νόμιμο πληκτρολόγιο και παραμένει σιωπηλό κατά την καραντίνα
        περνά όλους τους ελέγχους. Η μόνη άμυνα είναι η ερώτηση προς τον
        χρήστη: «Συνδέσατε μόλις ένα πληκτρολόγιο;»
  \item \textbf{Ο αναλυτής αποφασίζει.} Ο διαχωρισμός προνομίων περιορίζει τι
        μπορεί να αναγκάσει τον \en{root} να κάνει ένας παραβιασμένος αναλυτής,
        αλλά ο αναλυτής μπορεί ακόμη να εγκρίνει ή να απορρίψει συσκευές, ή να
        τερματίσει — κάτι που ανοίγει την πύλη. Η επιλογή αυτή είναι
        συνειδητή: η εναλλακτική (πύλη που μένει κλειστή όταν «πεθάνει» ο
        αναλυτής) θα μετέτρεπε κάθε τερματισμό από τον \en{OOM killer} σε
        κλείδωμα μηχανήματος με πληκτρολόγιο \en{USB}.
  \item \textbf{Οι θύρες διάσωσης.} Οι εσωτερικές θύρες και όσες δηλωθούν με
        \code{--allow-port} δεν ελέγχονται ούτε με κλειδωμένη οθόνη· όποιος
        έχει φυσική πρόσβαση μπορεί να τις χρησιμοποιήσει.
  \item \textbf{Οι κάρτες μνήμης δεν «φράζονται».} Μια κάρτα μέσα σε εγκεκριμένο
        αναγνώστη είναι μέσο, όχι συσκευή \en{USB}· το \code{--watch-media}
        ανιχνεύει, και το πολύ απενεργοποιεί ολόκληρο τον αναγνώστη.
  \item \textbf{Νέοι ριζικοί διανομείς.} Ελεγκτές που εμφανίζονται μετά την
        εκκίνηση (π.χ. σταθμός σύνδεσης) δεν κλείνονται.
  \item \textbf{Η αυτόματη προσάρτηση.} Κατά το στάδιο 4 η επιφάνεια εργασίας
        μπορεί να προσαρτήσει το μέσο. Ο κανόνας \en{udev} με
        \code{UDISKS\_IGNORE=1} το αποτρέπει μόνο για το \en{udisks}.
  \item \textbf{Πραγματικό υλικό επίθεσης.} Η αξιολόγηση έγινε με εξομοιωμένες
        συσκευές· η επικύρωση με \en{Rubber Ducky} και \en{O.MG} εκκρεμεί
        (§\ref{sec:eval-pending}).
\end{enumerate}

\section{Σχεδιαστικοί συμβιβασμοί}
\label{sec:tradeoffs}

\paragraph{Ασφάλεια έναντι διαθεσιμότητας.}
Σε ολόκληρη τη σχεδίαση, όταν οι δύο συγκρούονται, επιλέγεται η αποφυγή
κλειδώματος, εφόσον δεν οδηγεί σε σιωπηλή αποδοχή. Παράδειγμα: όταν ο πράκτορας
δεν είναι διαθέσιμος, μια νέα συσκευή \emph{κρατείται} και δεν απορρίπτεται, διότι
μια καταγεγραμμένη απόρριψη θα έκανε την επόμενη σύνδεση \sevC{}
(\code{previously-rejected}). Ομοίως, μια ερώτηση που έμεινε αναπάντητη
καταγράφεται ως «καμία απάντηση» και όχι ως απόρριψη.

\paragraph{Αυστηρότητα αποτυπωμάτων.}
Το ιστορικό χρησιμοποιεί \emph{κανονικοποιημένο} αποτύπωμα ώστε να μην
προκαλεί ψευδείς συναγερμούς μετάλλαξης όταν μια συσκευή αλλάζει θύρα
\en{USB}~2/3, ενώ ο χώρος εμπιστοσύνης χρησιμοποιεί \emph{ακατέργαστο}
αποτύπωμα. Το αποτέλεσμα είναι ότι μια απομνημονευμένη συσκευή μπορεί να
ερωτηθεί ξανά — μικρή ενόχληση, αλλά πάντα προς την ασφαλή κατεύθυνση.

\paragraph{\en{Python} έναντι \en{C}.}
Η \en{Python} εξαλείφει τα σφάλματα μνήμης στους αναλυτές μη έμπιστων
δεδομένων, με κόστος μεγαλύτερης έμπιστης βάσης (ο διερμηνέας) και αδυναμίας
χρονικού ορίου σε λειτουργίες πυρήνα που βρίσκονται σε αδιάκοπη αναμονή
(\en{uninterruptible sleep}). Το δεύτερο αντιμετωπίστηκε με διεργασίες-παιδιά.

\paragraph{Ανθρώπινος παράγοντας.}
Η τελική απόφαση ανήκει στον άνθρωπο, ο οποίος μπορεί να κάνει λάθος. Η
κλιμακούμενη τριβή στοχεύει να κάνει το λάθος δύσκολο εκεί όπου κοστίζει, χωρίς
να εκπαιδεύει τον χρήστη να πατά «Ναι» μηχανικά σε κάθε ερώτηση — φαινόμενο
γνωστό ως «κόπωση προειδοποιήσεων» (\en{warning fatigue}).

\section{Σύγκριση με το \en{USBGuard}}
\label{sec:vs-usbguard}

Τα δύο εργαλεία χρησιμοποιούν τον ίδιο μηχανισμό πυρήνα, αλλά απαντούν σε
διαφορετικό ερώτημα. Το \en{USBGuard} απαντά στο «επιτρέπεται αυτή η συσκευή
από την πολιτική;» και είναι κατάλληλο για διαχειριζόμενους στόλους
μηχανημάτων με γνωστό υλικό. Το \Probolos{} απαντά στο «είναι αυτή η άγνωστη
συσκευή αυτό που ισχυρίζεται;» και απευθύνεται σε χρήστη που συνδέει συχνά νέες
συσκευές. Τα δύο δεν είναι αμοιβαία αποκλειόμενα ως ιδέες, αλλά δεν πρέπει να
εκτελούνται ταυτόχρονα, διότι θα ανταγωνίζονταν για το ίδιο
\code{authorized\_default}.
